% Source: Corsin Nick/Class Notes/Week 9.pdf, pp. 2--8
\chapter{Week 9 --- Determinants, Volume \& the Length of a Curve}
\label{ch:week09}

% Transition: Claude Sonnet 5
This chapter asks what the determinant has to do with volume, and pushes the answer as far
as it goes: from the area of a parallelogram in $\mathbb{R}^2$, to the \newterm{Gram
determinant} of $m$ vectors in $\mathbb{R}^n$ for $m \neq n$, to the $d$-volume of a
parametrized $d$-submanifold of $\mathbb{R}^n$, with the length of a curve as the case
$d=1$. The exercise sheet is at the end of the chapter, in \cref{sec:week09_solutions}.

\section{Length, area, and volume}
\label{sec:length_area_volume}

\subsection{The area of a parallelogram in $\mathbb{R}^2$}
\label{sec:area_parallelogram}

The relationship between the determinant and volume is a key concept in this part of the
lecture. We illustrate the special case of $\mathbb{R}^2$.

\begin{proposition}[Area of a parallelogram]
\label{prop:area_parallelogram_r2}
Let $x,y \in \mathbb{R}^2$. The area of the parallelogram $(0,x,x+y,y)$ is
$$A = \big|\det(x\mid y)\big|.$$
\end{proposition}

\begin{proof}
Assume w.l.o.g.\ that $x = x_1e_1$ --- we justify this reduction below. Then, writing
$y = (y_1,y_2)\transp$,
$$A = |x_1y_2| = \left|\det\begin{pmatrix} x_1 & y_1 \\ 0 & y_2 \end{pmatrix}\right|,$$
directly from the base-times-height formula for a parallelogram with one side on the
$e_1$-axis.

% Sonnet 5 (Medium) --- FIG-W09-01
\begin{center}
\begin{tikzpicture}[>=Latex, scale=1.0]
  \draw[->] (-0.3,0) -- (3.6,0) node[right, font=\scriptsize] {$e_1$};
  \draw[->] (0,-0.3) -- (0,2.1) node[above, font=\scriptsize] {$e_2$};
  \coordinate (O) at (0,0);
  \coordinate (X) at (2.4,0);
  \coordinate (Y) at (0.7,1.7);
  \coordinate (XY) at (3.1,1.7);
  \draw[purple, dotted, thick] (O) -- (X) -- (XY) -- (Y) -- cycle;
  \draw[blue, thick, ->] (O) -- (X) node[below, font=\scriptsize] {$x=x_1e_1$};
  \draw[orange, thick, ->] (O) -- (Y) node[above left, font=\scriptsize] {$y$};
  \draw[green!50!black, thick, dotted] (Y) -- (0.7,0);
  \draw[green!50!black, ->] (0,-0.15) -- (0.7,-0.15) node[midway, below, font=\scriptsize] {$y_1$};
  \draw[green!50!black, ->] (0.85,0) -- (0.85,1.7) node[midway, right, font=\scriptsize] {$y_2$};
\end{tikzpicture}
\end{center}

For the reduction, let $R \in \operatorname{O}(2)$ be any rotation carrying $x,y$ to a pair
$Rx,Ry$ with $Rx$ along $e_1$; rotations preserve area, so
$$A = \big|\det(Rx\mid Ry)\big| = \big|\det(R)\det(x\mid y)\big| = \big|\det(x\mid y)\big|,$$
using $|\det R| = 1$ and multiplicativity of $\det$. So the \qt{w.l.o.g.} above was
justified: the formula $A=|\det(x\mid y)|$ holds for the rotated pair, and equals
$|\det(x\mid y)|$ for the original one.
\end{proof}

\subsection{The Gram determinant}
\label{sec:gram_determinant}

\begin{lemma}[Gram determinant of a square matrix]
\label{lem:gram_determinant_square}
For $M \in \mathrm{Mat}_{n\times n}(\mathbb{R})$,
$$|\det(M)| = \sqrt{\det(M)\det(M\transp)} = \sqrt{\det(M\transp)\det(M)}
  = \sqrt{\det(M\transp M)}.$$
\end{lemma}

\begin{proof}
The first equality is $|\det M| = \sqrt{(\det M)^2}$ together with $\det(M\transp)=\det(M)$;
the last is multiplicativity of the determinant, $\det(M\transp)\det(M)=\det(M\transp M)$.
\end{proof}

The quantity $\det(M\transp M)$ is called the \newterm{Gram determinant}
(\germanterm{Gram-Determinante}) of $M$. So, by \cref{prop:area_parallelogram_r2}, the
volume of an $n$-parallelotope in $\mathbb{R}^n$ is given by the Gram determinant, in
analogy to the determinant itself.

\begin{ainote}
The point of restating $|\det M|$ this way is that the middle expression, $\det(M\transp
M)$, still makes sense once $M \in \mathrm{Mat}_{n\times m}(\mathbb{R})$ is rectangular:
$M\transp M \in \mathrm{Mat}_{m\times m}(\mathbb{R})$ is square even when $M$ is not, so its
determinant is always defined, whereas $\det(M)$ itself is not. This is what lets the
notion of \qt{volume spanned by vectors} survive dropping the requirement $m=n$ --- a set of
fewer than $n$ vectors in $\mathbb{R}^n$ still spans a lower-dimensional parallelotope, and
that parallelotope still has an $m$-dimensional volume, even though no square matrix and no
ordinary determinant is available to compute it with.
\end{ainote}

\begin{definition}[Volume of a parallelotope via the Gram determinant]
\label{def:volume_gram_determinant}
For $m$ vectors $v_1,\dots,v_m \in \mathbb{R}^n$, the \newterm{area}/\newterm{length}/
\newterm{volume} of the parallelotope spanned by them is
$$\sqrt{\det\big(\langle v_i,v_j\rangle\big)} = \sqrt{\det(V\transp V)},
  \qquad V := (v_1\mid v_2\mid\cdots\mid v_m) \in \mathrm{Mat}_{n\times m}(\mathbb{R}).$$
\end{definition}

\begin{example}
\label{ex:gram_m1_length}
For $m=1$, $\sqrt{\det(v_1\transp v_1)} = \sqrt{\langle v_1,v_1\rangle} = \lVert v_1\rVert$
--- the length of $v_1$.
\end{example}

% Source: Corsin Nick/Class Notes/Week 9.pdf, p. 4
\begin{example}[$m=2,\ n=3$: area of a parallelogram in $\mathbb{R}^3$]
\label{ex:gram_m2_n3_area}
For $m=2$, $n=3$, we hope that the Gram determinant of
$$M = \begin{pmatrix} x_1 & y_1 \\ x_2 & y_2 \\ x_3 & y_3 \end{pmatrix}$$
gives the area of the parallelogram spanned by $x$ and $y$ in $\mathbb{R}^3$. If $\theta$ is
the angle between $x$ and $y$, this should be $A = \lVert x\rVert\lVert y\rVert|\sin\theta|$.

% Sonnet 5 (Medium) --- FIG-W09-02
\begin{center}
\begin{tikzpicture}[>=Latex, scale=1.3]
  \coordinate (O) at (0,0);
  \coordinate (X) at (2.6,1.7);
  \coordinate (Y) at (1.7,0.25);
  \coordinate (XY) at (4.3,1.95);
  \draw[orange, thick, ->] (O) -- (Y) node[below right, font=\scriptsize] {$y$};
  \draw[blue, thick, ->] (O) -- (X) node[above, font=\scriptsize] {$x$};
  \draw[purple, dotted, thick] (X) -- (XY) -- (Y);
  \draw[red, dotted, thick] (X) -- (2.6,0.36);
  \node[red, font=\scriptsize, right] at (2.68,1.0) {$\|x\|\sin\theta$};
  \draw (0.55,0.075) arc (5:33:0.6);
  \node[font=\scriptsize] at (0.9,0.32) {$\theta$};
\end{tikzpicture}
\end{center}

And indeed:
$$
\begin{aligned}
\sqrt{\det(M\transp M)}
  &= \sqrt{\det\begin{pmatrix} \langle x,x\rangle & \langle x,y\rangle \\
    \langle x,y\rangle & \langle y,y\rangle\end{pmatrix}} \\
  &= \sqrt{\langle x,x\rangle\langle y,y\rangle - \langle x,y\rangle^2} \\
  &= \sqrt{\lVert x\rVert^2\lVert y\rVert^2 - \lVert x\rVert^2\lVert y\rVert^2\cos^2\theta} \\
  &= \lVert x\rVert\lVert y\rVert|\sin\theta| \qquad (=\ \lvert x\times y\rvert).
\end{aligned}
$$
\end{example}

\subsection{Volume of embedded surfaces}
\label{sec:volume_embedded_surfaces}

% Source: Corsin Nick/Class Notes/Week 9.pdf, p. 5
% Def. 13.68
\begin{definition}[$d$-volume on a parametrized $d$-submanifold]
\label{def:d_volume_parametrized_submanifold}
Let $V \subseteq \mathbb{R}^d$ be open, and let $\phi : V \to \mathbb{R}^n$ be a parametrized
submanifold. Given a Jordan measurable set $E \subseteq V$ with $\overline E \subseteq V$,
we define the $d$-volume of $\phi(E)$ by
$$\vol_d(\phi(E)) := \int_E \sqrt{\det\big(D\phi(x)\transp D\phi(x)\big)}\,dx.$$
\end{definition}

Here the Gram determinant in the integrand is the volume of the parallelotope spanned by the
vectors $\partial_1\phi(x),\partial_2\phi(x),\dots,\partial_d\phi(x)$, which by assumption
(see \qt{parametrized submanifold}) are linearly independent --- i.e.\ $D\phi(x)$ has full
rank --- so the volume is strictly positive.

% Sonnet 5 (Medium) --- FIG-W09-03
\begin{center}
\begin{tikzpicture}[>=Latex, scale=1.3]
  \draw[thick] plot[smooth] coordinates {(-2.1,-0.4) (-0.8,0.65) (0.5,-0.15) (2.1,0.75)};
  \draw[thick] plot[smooth] coordinates {(-1.8,0.9) (-0.5,-0.5) (0.8,0.4) (2.3,-0.65)};
  \coordinate (P) at (0.0,0.05);
  \fill (P) circle (1.6pt) node[below left, font=\scriptsize] {$\phi(x)$};
  \draw[blue, thick, ->] (P) -- ++(1.2,0.4) coordinate (D1)
    node[right, font=\scriptsize] {$\partial_1\phi(x)$};
  \draw[purple, thick, ->] (P) -- ++(0.35,1.05) coordinate (D2)
    node[above, font=\scriptsize] {$\partial_2\phi(x)$};
  \draw[green!55!black, dashed]
    (D1) -- ++(0.35,1.05) -- (D2) -- cycle;
  \node[green!45!black, font=\scriptsize, align=center] at (2.1,1.5)
    {$\sqrt{D\phi(x)\transp D\phi(x)}$};
  \draw[green!55!black, thin, ->] (1.55,1.3) -- (0.95,1.0);
\end{tikzpicture}
\end{center}

This factor gives a measure for the stretching of $E$ as it is \qt{glued} to the surface. If
$$\sqrt{D\phi(x)\transp D\phi(x)} = 1 \quad \text{for all } x,$$
then $\phi$ is called an \newterm{isometry} (\germanterm{Isometrie}). A special case is the
length of a curve, treated in \cref{sec:length_of_curve} below.

\subsection{Example: area of a paraboloid patch}
\label{sec:example_paraboloid_area}

% Source: Corsin Nick/Class Notes/Week 9.pdf, pp. 6--7
\begin{example}[Area of a paraboloid patch]
\label{ex:paraboloid_patch_area}
Compute the area of the surface
$$F = \left\{(x,y,z)\in\mathbb{R}^3 \ :\ x^2+4y^2 < 8,\ \ z = \tfrac12(x^2+2y^2)\right\}.$$

Since $z = \tfrac12(x^2+2y^2)$ means $F$ is a graph, we can conveniently write $F$ as the
image of the parametrization
$$\phi : A \to F, \qquad (x,y) \mapsto \left(x,\,y,\,\tfrac12(x^2+2y^2)\right),
  \qquad A := \{(x,y)\in\mathbb{R}^2 : x^2+4y^2 < 8\}.$$

\textbf{Gram determinant.} We compute
$$D\phi(x,y) = \begin{pmatrix} 1 & 0 \\ 0 & 1 \\ x & 2y \end{pmatrix}
  \quad \Longrightarrow \quad
  D\phi(x,y)\transp D\phi(x,y) = \begin{pmatrix} 1+x^2 & 2xy \\ 2xy & 1+4y^2 \end{pmatrix},$$
so that
$$\sqrt{\det D\phi(x,y)\transp D\phi(x,y)}
  = \sqrt{(1+x^2)(1+4y^2) - 4x^2y^2}
  = \sqrt{1+x^2+4y^2}.$$

\textbf{The integral.} By \cref{def:d_volume_parametrized_submanifold},
$$\vol_2(F) = \int_A \sqrt{1+x^2+4y^2}\,dx\,dy, \qquad A = \{(x,y)\in\mathbb{R}^2 :
x^2+4y^2<8\}.$$

We should choose more convenient coordinates. Polar coordinates come to mind, but then we
still have problems, since $A$ is not radially symmetric. We want to write
$A = \Psi\big((0,R)\times(0,2\pi)\big)$ (up to a set of measure zero). Note that for $y=0$,
$|x| < \sqrt8 = 2\sqrt2$, and for $x=0$, $|y| < \sqrt2$. Choosing $R := \sqrt2$ and weighing
$x,y$ accordingly,
$$\begin{pmatrix} x \\ y \end{pmatrix} = \begin{pmatrix} 2r\cos\theta \\ r\sin\theta
\end{pmatrix} =: \Psi(r,\theta).$$

\textbf{Change of variables.}
$$dx\,dy = |\det D\Psi(r,\theta)|\,dr\,d\theta
  = \left\lVert \begin{matrix} 2\cos\theta & -2r\sin\theta \\ \sin\theta & r\cos\theta
    \end{matrix}\right\rVert dr\,d\theta = 2r\,dr\,d\theta,$$
giving
$$
\begin{aligned}
\vol_2(F) &= \int_A \sqrt{1+x^2+4y^2}\,dx\,dy
  = \int_0^{\sqrt2} dr\int_0^{2\pi} d\theta\ 2r\sqrt{1+4r^2} \\
  &= 4\pi\int_0^{\sqrt2} dr\ r(1+4r^2)^{1/2}
   = \frac\pi2\int_0^{\sqrt2} dr\ 8r(1+4r^2)^{1/2} \\
  &= \frac\pi2\left[(1+4r^2)^{3/2}\cdot\frac23\right]_0^{\sqrt2}
   = \frac\pi3\big(27-1\big) = \frac{26\pi}{3}.
\end{aligned}
$$
\end{example}

\begin{ainote}
Worth flagging what makes $A$ awkward: it is an ellipse, not a disc, so naive polar
coordinates $(x,y)=(\rho\cos\theta,\rho\sin\theta)$ would force $\rho$ to depend on $\theta$
along the boundary. The fix is to build the eccentricity into the change of variables itself
--- scaling $x$ by a factor of $2$ relative to $y$ turns the ellipse into the disc $r<\sqrt2$
before polar coordinates are even applied. This is the same idea as
\cref{ex:change_of_variables_domain} in Week 8: choose $\Psi$ to straighten the domain, not
to simplify the integrand.
\end{ainote}

\subsection{Length of a curve}
\label{sec:length_of_curve}

% Source: Corsin Nick/Class Notes/Week 9.pdf, p. 8
\begin{definition}[Length of a $C^1$-curve]
\label{def:length_c1_curve}
The \newterm{length} (\germanterm{Länge}) of a curve or path $\gamma \in
C^1([a,b],\mathbb{R}^n)$ is given by
$$L(\gamma) := \int_a^b \lvert\gamma'(t)\rvert\,dt.$$
\end{definition}

\begin{definition}[Length of a $C^0$-curve]
\label{def:length_c0_curve}
For a $C^0$-curve $c \in C^0([a,b],\mathbb{R}^n)$, the length is defined as
$$L(c) := \sup\left\{\sum_{i=1}^{n-1} \lvert c(t_{i+1})-c(t_i)\rvert
  \ :\ a = t_1 < \cdots < t_n = b\right\},$$
the supremum over all lengths of inscribed polygons.
\end{definition}

% Sonnet 5 (Medium) --- FIG-W09-04
\begin{center}
\begin{tikzpicture}[>=Latex, scale=1.0]
  \draw[thick] plot[smooth] coordinates
    {(0,0) (0.6,0.9) (1.3,0.5) (1.9,0.15) (2.5,0.75) (3.1,0.5) (3.7,1.3)};
  \node[font=\scriptsize] at (1.9,1.15) {$\operatorname{Im}(c)$};
  \draw[blue, thick]
    (0,0) -- (0.55,0.95) -- (1.35,0.55) -- (2.0,0.1) -- (2.55,0.7) -- (3.15,0.45) -- (3.7,1.25);
  \foreach \x/\y/\lbl/\pos in {0/0/{c(t_1)}/below,0.55/0.95/{c(t_2)}/above,1.35/0.55/{c(t_3)}/above,2.0/0.1/{c(t_4)}/below,2.55/0.7/{c(t_5)}/{above left},3.15/0.45/{c(t_6)}/below,3.7/1.25/{c(t_7)}/right} {
    \fill (\x,\y) circle (1.3pt);
    \node[font=\scriptsize, \pos] at (\x,\y) {$\lbl$};
  }
\end{tikzpicture}
\end{center}

\begin{remark}
For a $C^1$-curve, \cref{def:length_c1_curve,def:length_c0_curve} agree. Notice that
$$\lvert\gamma'(t)\rvert = \sqrt{D\gamma(t)\transp \cdot D\gamma(t)}$$
is exactly the Gram determinant of $D\gamma(t)$ in the sense of
\cref{def:volume_gram_determinant} --- the $m=1$ case of \cref{ex:gram_m1_length}, now
applied pointwise along the curve. This is why the length of a curve is the special case
$d=1$ of \cref{def:d_volume_parametrized_submanifold}: a curve is a $1$-dimensional
parametrized submanifold, and $L(\gamma) = \vol_1(\gamma([a,b]))$.
\end{remark}

\exercisesheet{9}

\textit{Statements quoted verbatim from \texttt{exercises/Ex9\_Analysis2\_eng.pdf} (assigned
20 April 2026, due 27 April 2026). Attribution: Prof.\ Joaquim Serra, D-MATH, ETH Z\"urich.}

% Source: Corsin Nick/Class Notes/Week 9.pdf, p. 1
\begin{ainote}[Corsin's recommendations]
Colour key, as given on the page: blue \textbf{= important}, orange \textbf{= semi-important},
red \textbf{= optional}. Only \textbf{9.1} and \textbf{9.6} carry an explicit colour marker
on the page; \textbf{9.2} is annotated with a hint but no clearly distinguishable priority
colour, so it is listed as \textbf{unmarked} below even though a hint is quoted for it.
\textbf{9.3--9.5} and \textbf{9.7} carry neither colour nor hint.

\begin{center}
\begin{tabular}{@{}l l p{0.52\linewidth}@{}}
\textbf{Problem} & \textbf{Priority} & \textbf{Corsin's note} \\
\hline
9.1 & important & (hint on coordinates --- see below) \\[3pt]
9.2 & unmarked & (hint on the domain of part 2 --- see below) \\[3pt]
9.3 & unmarked & --- \\[3pt]
9.4 & unmarked & --- \\[3pt]
9.5 & unmarked & --- \\[3pt]
9.6 & semi-important & ``This may become important in later courses. $\langle\cdot,\cdot\rangle$
  is the standard inner prod.'' \\[3pt]
9.7 & unmarked & --- \\
\end{tabular}
\end{center}
\end{ainote}

\begin{ainote}
Only \textbf{9.1} is tagged important above, but its companion problem \textbf{9.2} carries
an explicit worked hint from Corsin (unlike 9.3--9.5 and 9.7, which carry nothing at all), so
both are reproduced below with solutions. 9.3--9.7 are not reproduced.
\end{ainote}

% Source: exercises/Ex9_Analysis2_eng.pdf, p. 1
\begin{exercise}[9.1 --- Volume \textnormal{(important)}]
\label{ex:9.1}
Calculate the volume of the region $B \subset \mathbb{R}^3$ enclosed by the surfaces
$x^2+y^2+z^2=8$ and $2z=x^2+y^2$. \textit{Official hint:} use cylindrical coordinates.

% Source: Corsin Nick/Class Notes/Week 9.pdf, p. 1
\textit{Corsin's hint:}
$$\begin{pmatrix} x \\ y \\ z\end{pmatrix} = \begin{pmatrix} r\sin\theta\cos\varphi \\
  r\sin\theta\sin\varphi \\ r\cos\theta \end{pmatrix}
  \qquad \text{for } (r,\theta,\varphi) \in (0,\infty)\times(0,\pi)\times(0,2\pi).$$
\end{exercise}

\begin{ainote}
Corsin labels the formula above \qt{cylindrical coordinates}, but $(r\sin\theta\cos\varphi,\
r\sin\theta\sin\varphi,\ r\cos\theta)$ is the \textbf{spherical} parametrization, and his
range $\theta \in (-\pi/2,\pi/2)$ has been corrected to $(0,\pi)$, the range on which
$\theta \mapsto \cos\theta$ actually covers $[-1,1]$ bijectively., already
logged. The mislabelling is doubly worth noting here, since the official sheet's own hint
recommends \emph{cylindrical} coordinates $(r,\theta,z)$ --- and that is in fact the more
natural choice for this problem, as \cref{ex:solution_9.1} below uses: both surfaces are
surfaces of revolution about the $z$-axis, which cylindrical coordinates respect directly,
whereas spherical coordinates mix $r$ and $\theta$ into both defining equations.
\end{ainote}

% Source: exercises/Ex9_Analysis2_eng.pdf, p. 1
\begin{exercise}[9.2 --- Multiple integrals]
\label{ex:9.2}
\begin{enumerate}[label=\textbf{\arabic*.}]
  \item Let $D = [0,2]\times[0,1]$. Calculate $\displaystyle\iint_D (x^3+3x^2y+y^3)\,dx\,dy$.
  \item Let $D \subset \mathbb{R}^2$ be the interior of the triangle with vertices $(0,0)$,
    $(0,\pi)$, and $(\pi,\pi)$. Calculate $\displaystyle\iint_D x\cos(x+y)\,dx\,dy$.
  \item Let $D = \{(x,y)\in\mathbb{R}^2 \mid x>1,\ y>1,\ x+y<3\}$. Calculate
    $\displaystyle\iint_D \frac{1}{(x+y)^3}\,dx\,dy$.
\end{enumerate}

% Source: Corsin Nick/Class Notes/Week 9.pdf, p. 1
\textit{Corsin's hint (part 2):} you can parametrize the triangle as
$$D = \{(x,y)\in\mathbb{R}^2 : 0 \leq x \leq y \leq \pi\}.$$
\end{exercise}

\section{Solutions}
\label{sec:week09_solutions}

\begin{ainote}[TODO --- pending the remaining exercise statements]
Corsin's Week 9 notes contain no worked solutions (only the hints reproduced above), so all
solutions here are original. Only \textbf{9.1} and \textbf{9.2} are solved, matching the two
exercises reproduced above; 9.3--9.7 are left for a later pass.
\end{ainote}

\begin{exercisesolution}[Solution to \cref{ex:9.1}]
\label{ex:solution_9.1}
% Generator: Claude Sonnet 5 (Medium)
Following the official hint rather than Corsin's --- see the \texttt{ainote} above --- we use
cylindrical coordinates $(x,y,z) = (r\cos\theta,r\sin\theta,z)$, under which $x^2+y^2=r^2$.
The two bounding surfaces become $r^2+z^2=8$ and $z = r^2/2$; intersecting them,
$$z^2+2z-8 = 0 \quad\Longrightarrow\quad z = 2 \text{ or } z=-4,$$
and since $z=r^2/2\geq0$ on the paraboloid, the surfaces meet at $z=2$, $r=2$. For
$r \in (0,2)$ the paraboloid $z=r^2/2$ lies below the upper half of the sphere
$z=\sqrt{8-r^2}$, so
$$B = \left\{(r,\theta,z) : 0<r<2,\ 0<\theta<2\pi,\ \tfrac12r^2 < z < \sqrt{8-r^2}\right\}.$$
Hence, with $dx\,dy\,dz = r\,dr\,d\theta\,dz$,
$$
\begin{aligned}
\vol(B) &= \int_0^{2\pi}\!\!\int_0^2 r\left(\sqrt{8-r^2} - \frac{r^2}{2}\right)dr\,d\theta
  = 2\pi\int_0^2 \left(r\sqrt{8-r^2} - \frac{r^3}{2}\right)dr.
\end{aligned}
$$
For the first term, substitute $u = 8-r^2$, $du=-2r\,dr$:
$$\int_0^2 r\sqrt{8-r^2}\,dr = \left[-\frac{u^{3/2}}{3}\right]_{u=8}^{u=4}
  = \frac{8^{3/2}-4^{3/2}}{3} = \frac{16\sqrt2-8}{3}.$$
For the second, $\int_0^2 \tfrac12 r^3\,dr = \tfrac12\cdot\tfrac{16}{4} = 2$. So
$$\vol(B) = 2\pi\left(\frac{16\sqrt2-8}{3} - 2\right) = 2\pi\cdot\frac{16\sqrt2-14}{3}
  = \boxed{\ \frac{4\pi}{3}\big(8\sqrt2-7\big)\ }.$$
\end{exercisesolution}

\begin{exercisesolution}[Solution to \cref{ex:9.2}]
% Generator: Claude Sonnet 5 (Medium)
\textbf{Part 1.} By Fubini (\cref{thm:fubini}),
$$\iint_D (x^3+3x^2y+y^3)\,dx\,dy = \int_0^2\int_0^1 (x^3+3x^2y+y^3)\,dy\,dx
  = \int_0^2\left(x^3+\frac32x^2+\frac14\right)dx.$$
The last integral is $\left[\tfrac{x^4}{4}+\tfrac{x^3}{2}+\tfrac{x}{4}\right]_0^2 = 4+4+\tfrac12
= \boxed{\ \tfrac{17}{2}\ }$.

\textbf{Part 2.} Using Corsin's parametrization $D=\{0\leq x\leq y\leq\pi\}$,
$$\iint_D x\cos(x+y)\,dx\,dy = \int_0^\pi\int_x^\pi x\cos(x+y)\,dy\,dx
  = \int_0^\pi x\big[\sin(x+y)\big]_{y=x}^{y=\pi}\,dx.$$
Since $\sin(x+\pi) = -\sin x$, the bracket equals $-\sin x - \sin(2x)$, so
$$\iint_D x\cos(x+y)\,dx\,dy = -\int_0^\pi x\sin x\,dx - \int_0^\pi x\sin(2x)\,dx.$$
Integrating by parts, $\int_0^\pi x\sin x\,dx = [-x\cos x+\sin x]_0^\pi = \pi$, and
$\int_0^\pi x\sin(2x)\,dx = \left[-\tfrac{x}{2}\cos(2x)+\tfrac14\sin(2x)\right]_0^\pi
= -\tfrac\pi2$. Hence
$$\iint_D x\cos(x+y)\,dx\,dy = -\pi - \left(-\frac\pi2\right) = \boxed{\ -\frac\pi2\ }.$$

\textbf{Part 3.} Here $D = \{1<x<2,\ 1<y<3-x\}$, so
$$\iint_D \frac{1}{(x+y)^3}\,dx\,dy = \int_1^2\int_1^{3-x} (x+y)^{-3}\,dy\,dx
  = \int_1^2 \left[-\frac12(x+y)^{-2}\right]_{y=1}^{y=3-x}dx.$$
Since $x+(3-x)=3$, the bracket is $-\tfrac12\big(3^{-2}-(x+1)^{-2}\big)
= \tfrac12(x+1)^{-2} - \tfrac1{18}$, and
$$\int_1^2\left(\frac{1}{2(x+1)^2}-\frac1{18}\right)dx
  = \left[-\frac{1}{2(x+1)}\right]_1^2 - \frac1{18}
  = \left(-\frac16+\frac14\right) - \frac1{18} = \frac1{12}-\frac1{18}
  = \boxed{\ \frac1{36}\ }.$$
\end{exercisesolution}

\begin{ainote}
All three parts of \cref{ex:9.2} and \cref{ex:9.1} were worked out independently before
consulting \texttt{Sol9\_Analysis2\_eng.pdf}; all four answers ($\tfrac{4\pi}{3}(8\sqrt2-7)$,
$\tfrac{17}{2}$, $-\tfrac\pi2$, $\tfrac1{36}$) match the official solution exactly, so no
divergence to record.
\end{ainote}
